%% kapa: exam=1 round=36 year=2025 period=1 form=A subject=civil_law subject_ko=민법 seq=31 number=31
%% source: Q-Net 감정평가사 기출문제 (cst00309) artlSeq=5231525 file=raw/1008/5231525/2025년 제36회 감정평가사 1차시험 1교시 문제지 원본.pdf
%% conversion: kordoc 4.24.51 -> tools/gen1.py (auto); review=proofread-sonnet; flags=none
\begin{question}{31}
지상권에 관한 설명으로 옳지 않은 것은? (다툼이 있으면 판례에 따름)
\begin{choices}
  \item 지상권에 저당권을 설정해 준 지상권자가 지상권의 목적인 토지를 매수한 때에는 지상권이 혼동으로 소멸하지 않는다.
  \item 토지에 저당권, 지상권, 저당권이 순차적으로 설정된 경우, 나중에 설정된 저당권이 실행되면 지상권은 소멸한다.
  \item 지상권자는 제3자에게 구분지상권을 설정해 줄 수 있다.
  \item 토지의 담보가치 하락을 막기 위해 설정된 지상권은 피담보채권이 소멸하면 존속기간과 관계없이 소멸한다.
  \item 지상권에 저당권이 설정된 경우, 지상권설정자의 지상권소멸청구는 그 저당권자에게 통지한 후 상당한 기간이 경과해야 효력이 생긴다.
\end{choices}
\end{question}
\answer{3}
